\answer{ex:22:1}{(а)~$32\%$ і $100\%$. (б)~$R_\uparrow-R_\downarrow\ge2\sqrt{2000\,\text{Гц}/1\,\text{с}}\approx89$~Гц, лише $4.5\%$ сумарної частоти.}
\answer{ex:22:2}{(а)~$n\ge4(p_1q_1+p_2q_2)/\Delta^2\approx560$. (б)~$\approx56$.}
\answer{ex:22:3}{(а)~Детальний баланс: $\rho PK$ симетричне. (б)~Частка промахів дорівнює середньому гармонічному $1/\langle1/P\rangle=0.18$ проти $0.5$ без зворотного зв'язку.}
\answer{ex:22:4}{(а)~Паралелограм площею $4h^2=0.04$ (з урахуванням обмеження $p$ --- чисельно $0.045$). (б)~$\approx0.18$.}
\answer{ex:22:5}{$0.49$: частина культур іде в область, де промах майже завжди, і блукає там. З обмеженням --- $1.00$: на обмеженій множині густина $\propto1/P$ збирається в поглинальній області.}
\answer{ex:22:6}{(а)~Потрібно $\ge5$ рівнів; $8$ електродів вистачає. (б)~Ні: потрібно $r_{\max}\ge25/T=100$~Гц.}
\answer{ex:22:7}{Відкрита задача. Час випадкового пошуку зростає з $d$ приблизно як відношення об'ємів $(L/h)^d$, у пошуку зі знаком помилки --- лінійно за $d$. Розділяє гіпотези дослід із перестановкою: після промаху --- передбачуваний, але сильний стимул, після удару --- непередбачуваний, але слабкий; потрібно кілька десятків культур у кожній групі.}
